% =====================================================================
%  CHAPTER 2: 医学模式
% =====================================================================
\chapter{医学模式}
\setcounter{chapter}{2}
\chapterintro{
掌握医学模式的概念及各种医学模式的基本观点；重点掌握生物-心理-社会医学模式的概念、基本内涵及其意义（\textbf{教学难点}）；理解健康观的演变和健康与疾病概念的扩展。
}

% ----- 2.1 医学模式的概念 -----
\section{医学模式的概念}

\begin{defbox}
\textbf{医学模式（Medical Model）}：指人们在观察和处理健康与疾病问题的实践中形成的\keyword{总体认识框架和基本观点}。它反映了特定历史时期医学科学发展的整体水平、医学思维方式和对健康与疾病问题的基本认知方式，是人类与疾病作斗争、获取健康经验的总结。

医学模式是\imp{医学实践活动的指导思想}，影响和决定着医学教育、医学研究、临床诊疗和卫生服务的方向与模式。
\end{defbox}

\begin{tipbox}
医学模式的核心内容包括三个方面：
\begin{enumerate}[leftmargin=*]
    \item 疾病观——对疾病本质、原因和机制的基本认识
    \item 健康观——对健康本质和标准的认识
    \item 医学观——对医学目的、任务和方法的总体看法
\end{enumerate}
\end{tipbox}

% ----- 2.2 医学模式的演变 -----
\section{医学模式的演变过程}

\subsection{神灵主义医学模式}

\begin{defbox}
\textbf{神灵主义医学模式（Spiritualistic Medical Model）}：原始社会中，人们将健康与疾病归因于\keyword{超自然力量}（神灵、鬼怪等），认为疾病是神灵惩罚或邪魔附体所致，健康是神灵赐福。对疾病的治疗主要依靠巫术、祈祷等驱鬼降神手段。
\end{defbox}

\textbf{基本观点：}
\begin{itemize}[leftmargin=*]
    \item 疾病 = 超自然力量的惩罚
    \item 治疗 = 巫术、祈祷、驱魔
    \item 医巫不分，健康与疾病被神秘化
\end{itemize}

\subsection{自然哲学医学模式}

\begin{defbox}
\textbf{自然哲学医学模式（Natural Philosophical Medical Model）}：古代文明时期，人们开始用\keyword{自然哲学思想}来解释健康与疾病现象，将人体视为自然界的缩影，强调人与环境的和谐统一。
\end{defbox}

\textbf{基本观点与代表：}
\begin{enumerate}[leftmargin=*]
    \item \imp{古希腊希波克拉底}——体液学说（四体液：血液、黏液、黄胆汁、黑胆汁），四种体液平衡即为健康。强调环境和生活方式对健康的影响，《空气、水与地域》阐述了环境与疾病的关系
    \item \imp{中国《黄帝内经》}——阴阳五行学说、天人合一思想。强调情志、饮食、起居、劳逸等对健康的影响，提出"上工治未病"的预防思想
    \item \imp{古印度}——三体液学说（气、胆、痰）
\end{enumerate}

\begin{keybox}
自然哲学医学模式的核心特点：
\begin{itemize}[leftmargin=*]
    \item \keyword{整体观}：人体是一个整体，人与自然也是一个整体
    \item \keyword{平衡观}：健康是体内各要素的平衡，疾病是失衡
    \item 从自然哲学角度解释疾病，摆脱了纯粹的神灵迷信
    \item 开始重视环境、生活方式等因素对健康的影响
\end{itemize}
\end{keybox}

\subsection{机械论医学模式}

\begin{defbox}
\textbf{机械论医学模式（Mechanistic Medical Model）}：文艺复兴后，随着物理学和机械力学的兴起，人们将人体视为一部\keyword{精密的机器}，疾病是机器零件损坏或功能失常，医生的任务是修理机器。
\end{defbox}

\textbf{基本观点：}
\begin{itemize}[leftmargin=*]
    \item 人体 = 精密的机器（心脏=泵、关节=杠杆、肺=风箱）
    \item 疾病 = 机器部件损坏
    \item 治疗 = 修复损坏的部件
    \item 代表：笛卡尔（"动物是机器"）、拉美特利（"人是机器"）
\end{itemize}

\textbf{历史贡献与局限：}
\begin{itemize}[leftmargin=*]
    \item 贡献：推动了实验医学和解剖学的发展，使医学摆脱了思辨哲学的束缚
    \item 局限：忽视了人体生命的复杂性和整体性；忽视了心理、社会因素的作用；将复杂的生命活动简化为机械运动
\end{itemize}

\subsection{生物医学模式}

\begin{defbox}
\textbf{生物医学模式（Biomedical Model）}：以\keyword{生物学为基础}，从生物学角度认识和解释健康与疾病的医学模式。认为任何疾病都可以在器官、组织、细胞或分子水平上找到生物学的异常改变，每种疾病都可以找到特定的生物学病因和发病机制。
\end{defbox}

\begin{keybox}
\textbf{生物医学模式的基本观点（理论基础）：}
\begin{enumerate}[leftmargin=*]
    \item \imp{病因学一元论}：每种疾病都有特定的生物或理化致病因子（病原体、毒素、理化因素等）
    \item \imp{器官定位论}：疾病定位于特定的器官、组织或细胞
    \item \imp{特异性治疗论}：针对特定病因采用特异性治疗方法
    \item \imp{还原论}：复杂的生命现象可以还原为物理、化学过程来解释
    \item \imp{心身二元论}：将躯体与精神分离，忽视心理因素的致病作用
\end{enumerate}
\end{keybox}

\begin{tipbox}
\textbf{生物医学模式的历史贡献：}
\begin{enumerate}[leftmargin=*]
    \item 推动了现代医学的巨大进步——微生物学（科赫法则）、免疫学、麻醉、消毒技术、抗生素等的发现和应用
    \item 在传染病和急性病的防治中取得了辉煌成就，显著降低死亡率和发病率
    \item 建立了科学的疾病分类体系和诊断治疗方法
    \item 促进了医学的专科化和精细化发展
\end{enumerate}

\textbf{生物医学模式的局限性：}
\begin{enumerate}[leftmargin=*]
    \item 忽视了心理、社会因素在疾病发生发展中的作用
    \item 将人简单还原为生物学个体，忽视了人的社会属性
    \item 无法解释现代疾病谱的变化——慢性非传染性疾病、心身疾病等
    \item 过度依赖技术和药物，导致医疗费用不断攀升
    \item 忽视预防和健康促进的作用
    \item 无法解释一些功能性疾病的病因
\end{enumerate}
\end{tipbox}

% ----- 2.3 生物-心理-社会医学模式 -----
\section{生物-心理-社会医学模式}

\subsection{产生的背景}

\begin{keybox}
\textbf{生物-心理-社会医学模式产生的四大背景：}
\begin{enumerate}[leftmargin=*]
    \item \imp{疾病谱和死因谱的转变}——从传染病为主转向慢性非传染性疾病（心脑血管病、恶性肿瘤、糖尿病等）为主。这些疾病的病因中，行为生活方式、环境、社会因素占重要地位
    \item \imp{对健康需求的提高}——人们不再满足于不生病，而是追求身体、心理和社会的完好状态
    \item \imp{医学科学的发展}——心理学、社会学、人类学等与医学的交叉融合，心身医学、行为医学的发展
    \item \imp{健康社会决定因素的认知}——大量研究证据表明社会因素对健康的影响至关重要
\end{enumerate}
\end{keybox}

\textbf{医学模式转变的流行病学证据：}
\begin{itemize}[leftmargin=*]
    \item 美国：前10位死因中，行为生活方式因素占50\%以上，环境因素占20\%，生物学因素占20\%，卫生服务仅占10\%
    \item 中国：随着经济发展和老龄化，慢性病负担日益加重，危险因素以行为生活方式和社会环境因素为主
\end{itemize}

\subsection{恩格尔的生物-心理-社会医学模式}

\begin{defbox}
\textbf{Engel的生物-心理-社会医学模式（Bio-Psycho-Social Medical Model, 1977）}：由美国罗切斯特大学精神病学和内科学教授G.L. Engel在《科学》杂志上发表论文正式提出。该模式认为，\keyword{健康与疾病是生物、心理和社会因素相互作用的结果}，要全面认识健康和疾病，必须考虑生物、心理和社会三个层面的因素及其相互作用。
\end{defbox}

\begin{keybox}
\imp{（教学难点*）}\textbf{生物-心理-社会医学模式的基本内涵：}
\begin{enumerate}[leftmargin=*]
    \item \imp{多层次系统论}：人体是一个多层次的等级系统，从分子→细胞→组织→器官→系统→个体→家庭→社区→社会，每一层次既有其自身的特点和规律，又受上一层和下一层的影响。健康与疾病是各层次因素相互作用的结果。

    \item \imp{三维健康决定论}：健康与疾病同时受生物因素（遗传、病原体、生理功能等）、心理因素（认知、情绪、人格、压力等）和社会因素（经济、文化、制度、环境等）三个维度的影响，三个维度相互交织、互为因果。

    \item \imp{整体论}：人是一个完整的系统，不能还原为各个部分的简单叠加。医生需要将病人视为一个"完整的人"，而非仅仅是某种疾病的载体。

    \item \imp{患者中心论}：医疗服务应以患者为中心，关注患者的主观体验、心理状态和社会处境，而非仅关注疾病的客观指标。
\end{enumerate}
\end{keybox}

\begin{defbox}
\textbf{环境健康医学模式（Blum's Model, 1974）}：由Blum提出，将影响健康的因素归纳为四大类：\keyword{环境因素、行为生活方式、生物遗传因素和医疗卫生服务}。其中\imp{环境因素是最重要的健康决定因素}。

\textbf{综合健康医学模式（Lalonde-Dever模型）}：在Blum模型基础上的进一步发展，将影响健康的因素分为：①人类生物学；②环境；③生活方式；④卫生保健体系。
\end{defbox}

\subsection{生物-心理-社会医学模式的意义}

\begin{keybox}
\textbf{生物-心理-社会医学模式的实践意义：}
\begin{enumerate}[leftmargin=*]
    \item \imp{更新医学观念}——促进健康观的更新，从"没有疾病"扩展到"躯体、心理和社会适应的完好状态"
    \item \imp{推动医学教育改革}——医学教育从单纯生物学向生物-心理-社会综合知识体系转变，加强人文社会科学教育
    \item \imp{改变医疗服务模式}——从"以疾病为中心"转向"以患者为中心"，重视医患沟通、心理关怀和社会支持
    \item \imp{促进预防医学发展}——从"三级预防"发展为更加重视"零级预防"和社会预防
    \item \imp{推动卫生服务体系改革}——加强初级卫生保健、社区卫生服务和全科医学发展
    \item \imp{促进跨学科合作}——医学与心理学、社会学、经济学、管理学等学科的交叉融合
    \item \imp{指导卫生政策制定}——强调多部门协作和社会参与，影响"将健康融入所有政策"的理念
\end{enumerate}
\end{keybox}

% ----- 2.4 健康观的演变 -----
\section{健康观的演变}

\subsection{消极健康观}

\begin{defbox}
\textbf{消极健康观（经典健康观）}：将健康定义为\keyword{"没有疾病"}（Health is the absence of disease）。这种观点以疾病的生物学异常为参照，健康仅仅被理解为非疾病状态。
\end{defbox}

\textbf{消极健康观的特点与局限：}
\begin{itemize}[leftmargin=*]
    \item 以疾病为中心定义健康，标准明确但过于狭窄
    \item 忽视了健康的多维度性（心理、社会维度）
    \item 无法涵盖亚健康状态和功能健康
    \item 与"生物医学模式"相适应
\end{itemize}

\subsection{WHO健康观}

\begin{keybox}
\textbf{世界卫生组织（WHO, 1948）的健康定义：}\\
\textit{"Health is a state of complete physical, mental and social well-being, and not merely the absence of disease or infirmity."}\\
\textbf{健康不仅是躯体没有疾病或虚弱，而且是身体、心理和社会适应的完好状态。}
\end{keybox}

\textbf{WHO健康定义的深远意义：}
\begin{enumerate}[leftmargin=*]
    \item 首次提出\imp{三维健康观}——突破了"没有疾病=健康"的传统观念
    \item 强调\imp{社会维度}——社会适应良好是健康的重要组成部分
    \item 将健康从\imp{"消极"定义转向"积极"定义}
    \item 蕴含\imp{积极健康观}的理念——健康是积极追求的理想状态
    \item 为新医学模式（生物-心理-社会医学模式）提供了理论基础
    \item 推动医疗服务从"治病"转向"促进健康"
\end{enumerate}

\subsection{健康与疾病概念的扩展}

\begin{defbox}
\textbf{疾病（Disease）、病患（Illness）与患病（Sickness）的区分：}
\begin{itemize}[leftmargin=*]
    \item \imp{Disease（疾病）}：医学上客观存在的病理生理学异常，可通过检查诊断确认，是\keyword{生物学层面的异常}
    \item \imp{Illness（病患）}：患者主观感受到的不适、痛苦和功能受限等体验，是\keyword{心理学层面的感受}
    \item \imp{Sickness（患病）}：社会对个体疾病状态的认可和反应，涉及社会角色、权利和义务，是\keyword{社会学层面的状态}
\end{itemize}
\end{defbox}

\begin{tcolorbox}[colback=light, colframe=main, boxrule=0.5pt, arc=4pt, title=\textbf{现代健康观的多维视角}, breakable]
\begin{enumerate}[leftmargin=*]
    \item \imp{躯体健康}：身体结构完整、生理功能正常
    \item \imp{心理健康}：认知功能正常、情绪稳定、人格健全、良好的应对能力
    \item \imp{社会适应良好}：能有效扮演社会角色、维持良好的人际关系和社会功能
    \item \imp{道德健康}：行为符合社会道德规范，具有社会责任感和道德判断力
    \item \imp{亚健康状态}：介于健康与疾病之间的中间状态，表现为功能减退但无明显器质性病变
\end{enumerate}
\end{tcolorbox}

% ----- 复习题 -----
\section{复习题}

\subsection*{一、选择题}
\begin{enumerate}[leftmargin=*, itemsep=3pt]
    \item 生物-心理-社会医学模式的提出者是：
    \begin{enumerate}
        \item[(A)] Virchow
        \item[(B)] Blum
        \item[(C)] Engel
        \item[(D)] Lalonde
    \end{enumerate}
    \item 以下哪种医学模式将人体视为"机器"？
    \begin{enumerate}
        \item[(A)] 神灵主义医学模式
        \item[(B)] 自然哲学医学模式
        \item[(C)] 机械论医学模式
        \item[(D)] 生物医学模式
    \end{enumerate}
    \item WHO健康定义的时间是：
    \begin{enumerate}
        \item[(A)] 1946年
        \item[(B)] 1948年
        \item[(C)] 1977年
        \item[(D)] 1978年
    \end{enumerate}
    \item 影响健康的环境因素中，Blum模型认为最重要因素是：
    \begin{enumerate}
        \item[(A)] 生物遗传因素
        \item[(B)] 行为生活方式
        \item[(C)] 环境因素
        \item[(D)] 医疗卫生服务
    \end{enumerate}
    \item 以下属于Disease（疾病）定义的是：
    \begin{enumerate}
        \item[(A)] 患者主观感受的不适
        \item[(B)] 客观的病理生理学异常
        \item[(C)] 社会对疾病状态的认可
        \item[(D)] 社会角色功能障碍
    \end{enumerate}
\end{enumerate}

\subsection*{二、名词解释}
\begin{enumerate}[leftmargin=*, itemsep=3pt]
    \item \textbf{医学模式（Medical Model）}
    \item \textbf{生物-心理-社会医学模式}
    \item \textbf{生物医学模式}
    \item \textbf{健康（WHO, 1948定义）}
\end{enumerate}

\subsection*{三、简答题}
\begin{enumerate}[leftmargin=*, itemsep=3pt]
    \item 简述生物-心理-社会医学模式产生的背景。
    \item 试述生物-心理-社会医学模式的基本内涵及其意义。
    \item 简述生物医学模式的基本观点及其局限性。
    \item 试述Disease、Illness、Sickness三者的区别与联系。
    \item 简述WHO健康定义的内涵和意义。
    \item 试论医学模式转变对医疗卫生实践的影响。
\end{enumerate}

\subsection*{参考答案要点}

\subsubsection*{一、选择题}
1. (C)；2. (C)；3. (B)；4. (C)；5. (B)

\subsubsection*{二、名词解释}
\begin{enumerate}[leftmargin=*, itemsep=3pt]
    \item \textbf{医学模式}：人们在观察和处理健康与疾病问题实践中形成的总体认识框架和基本观点，是医学实践活动的指导思想。
    \item \textbf{生物-心理-社会医学模式}：Engel于1977年提出，认为健康与疾病是生物、心理和社会因素相互作用的结果，需从三个层面全面认识健康和疾病。
    \item \textbf{生物医学模式}：以生物学为基础，认为任何疾病都可在器官、组织、细胞或分子水平找到生物学异常改变的医学模式。
    \item \textbf{健康（WHO）}：健康不仅是躯体没有疾病或虚弱，而且是身体、心理和社会适应的完好状态（1948年WHO宪章）。
\end{enumerate}

\subsubsection*{三、简答题}
\begin{enumerate}[leftmargin=*]
    \item 四大背景：①疾病谱和死因谱转变（慢性病为主）；②健康需求提高（从"不生病"到"完好状态"）；③医学科学发展（交叉融合）；④健康社会决定因素认知（证据积累）。
    \item 基本内涵：多层次系统论；三维健康决定论；整体论；患者中心论。意义：更新医学观念、推动教育改革、改变服务模式、促进预防、推动体系改革等。
    \item 五大基本观点：病因学一元论、器官定位论、特异性治疗论、还原论、心身二元论。局限：忽视心理社会因素、无法解释疾病谱变化、过度依赖技术。
    \item Disease是客观病理学异常（生物学层面），Illness是患者主观感受（心理学层面），Sickness是社会角色状态（社会学层面）。三者既有区别又相互关联。
    \item WHO定义包括躯体、心理、社会三维完好状态。意义：突破传统消极健康观，为生物-心理-社会医学模式提供理论基础。
    \item 从生物医学向生物-心理-社会医学的转变体现在：观念更新、教育变革、服务模式改变、预防策略调整等方面。
\end{enumerate}

\newpage
